% GENERATED FILE. Edit canonical lessons, then run npm run generate:books.
% canonical-source-hash: fnv1a64:da6887ddec51ba27
% canonical-lessons: TE-C335-kiritam, TE-C335-simhasanam, TE-C335-pattabhisekam, TE-C335-pattam-kattu, TE-C335-palana

\chapter{A Crown for the King}
\label{ch:te-a-crown-for-the-king}
\hlchaptermodality{335}

\begin{chapteropening}
\textbf{By the end of this chapter:} \emph{I can say a crown, a throne, a coronation, to enthrone and a reign.}

Complete the last lesson of chapter 335: I can say a crown, a throne, a coronation, to enthrone and a reign.
\end{chapteropening}

\section[\texorpdfstring{\te{కిరీటం} (kirīṭaṁ)}{కిరీటం (kirīṭaṁ)}]{\te{కిరీటం} (kirīṭaṁ) — a crown}

\label{lesson:TE-C335-kiritam}



\begin{quote}
Before the new one: say the Telugu for to take birth, then the Telugu for to bless.
\end{quote}

\subsection*{You'll want to know: \te{కిరీటం}}
\textbf{\te{కిరీటం}} — \emph{kirīṭaṁ} — \textquotedblleft{}a crown\textquotedblright{}.

\begin{quote}
\te{రాజు} \te{తల} \te{మీద} \te{బంగారు} \te{కిరీటం} \te{ఉంది}.
\end{quote}

\begin{quote}
\emph{rāju tala mīda baṅgāru kirīṭaṁ undi.}
\end{quote}

\begin{quote}
There is a gold crown on the king's head.
\end{quote}

Plural: \te{కిరీటాలు}, \emph{kirīṭālu}.

A Sanskrit -a noun in -\te{ం}; its plural is \te{కిరీటాలు}.

\begin{cousinweb}
\textbf{\te{కిరీటం}} (\emph{kirīṭaṁ}) is borrowed from Sanskrit \emph{kirīṭa}, \textquotedblleft{}a diadem, a crest\textquotedblright{} (CDIAL 3176).
\end{cousinweb}

\begin{grammarlens}[title={What you've built}]
One more word for telling what happened.
\end{grammarlens}

\subsection*{Your turn}
\begin{itemize}
\raggedright
  \item \emph{Say it:} \emph{kirīṭaṁ}
  \item \emph{Say it:} \emph{kirīṭaṁ}, once more
  \item \emph{Say it:} \emph{rāju tala mīda baṅgāru kirīṭaṁ undi.} — There is a gold crown on the king's head.
  \item \emph{Recall:} say the Telugu for to take birth, then the Telugu for to bless, then say \emph{kirīṭaṁ} again
\end{itemize}

\begin{tcolorbox}[breakable,colback=teal!4,colframe=teal!35!black,title={Before you move on}]
What does \te{కిరీటం} mean? (\textquotedblleft{}A crown\textquotedblright{}.) Say it once more.
\end{tcolorbox}
\section[\texorpdfstring{\te{సింహాసనం} (siṁhāsanaṁ)}{సింహాసనం (siṁhāsanaṁ)}]{\te{సింహాసనం} (siṁhāsanaṁ) — a throne}

\label{lesson:TE-C335-simhasanam}



\begin{quote}
Before the new one: say the Telugu for to bless, then the Telugu for a crown.
\end{quote}

\subsection*{You'll want to know: \te{సింహాసనం}}
\textbf{\te{సింహాసనం}} — \emph{siṁhāsanaṁ} — \textquotedblleft{}a throne\textquotedblright{}.

\begin{quote}
\te{రాజు} \te{సింహాసనం} \te{మీద} \te{కూర్చున్నాడు}.
\end{quote}

\begin{quote}
\emph{rāju siṁhāsanaṁ mīda kūrcunnāḍu.}
\end{quote}

\begin{quote}
The king sat on the throne.
\end{quote}

Literally a lion's seat: \te{సింహం}, a lion, and \te{ఆసనం}, a seat.

\begin{cousinweb}
\textbf{\te{సింహాసనం}} (\emph{siṁhāsanaṁ}) is borrowed from Sanskrit \emph{siṁhāsana}, \textquotedblleft{}the lion's seat, a throne\textquotedblright{} (Monier-Williams), from \emph{siṁha}, \textquotedblleft{}a lion\textquotedblright{}, and \emph{āsana}, \textquotedblleft{}a seat\textquotedblright{}.
\end{cousinweb}

\begin{grammarlens}[title={What you've built}]
One more word for telling what happened.
\end{grammarlens}

\subsection*{Your turn}
\begin{itemize}
\raggedright
  \item \emph{Say it:} \emph{siṁhāsanaṁ}
  \item \emph{Say it:} \emph{siṁhāsanaṁ}, once more
  \item \emph{Say it:} \emph{rāju siṁhāsanaṁ mīda kūrcunnāḍu.} — The king sat on the throne.
  \item \emph{Recall:} say the Telugu for to bless, then the Telugu for a crown, then say \emph{siṁhāsanaṁ} again
\end{itemize}

\begin{tcolorbox}[breakable,colback=teal!4,colframe=teal!35!black,title={Before you move on}]
What does \te{సింహాసనం} mean? (\textquotedblleft{}A throne\textquotedblright{}.) Say it once more.
\end{tcolorbox}
\section[\texorpdfstring{\te{పట్టాభిషేకం} (paṭṭābhiṣēkaṁ)}{పట్టాభిషేకం (paṭṭābhiṣēkaṁ)}]{\te{పట్టాభిషేకం} (paṭṭābhiṣēkaṁ) — a coronation}

\label{lesson:TE-C335-pattabhisekam}



\begin{quote}
Before the new one: say the Telugu for a crown, then the Telugu for a throne.
\end{quote}

\subsection*{You'll want to know: \te{పట్టాభిషేకం}}
\textbf{\te{పట్టాభిషేకం}} — \emph{paṭṭābhiṣēkaṁ} — \textquotedblleft{}a coronation\textquotedblright{}.

\begin{quote}
\te{కొత్త} \te{రాజు} \te{పట్టాభిషేకం} \te{నిన్న} \te{జరిగింది}.
\end{quote}

\begin{quote}
\emph{kotta rāju paṭṭābhiṣēkaṁ ninna jarigindi.}
\end{quote}

\begin{quote}
The new king's coronation took place yesterday.
\end{quote}

\te{పట్టం}, the royal band, and \te{అభిషేకం}, an anointing with water.

\begin{cousinweb}
\textbf{\te{పట్టాభిషేకం}} (\emph{paṭṭābhiṣēkaṁ}) is borrowed from Sanskrit: \emph{paṭṭa}, \textquotedblleft{}a band, a diadem\textquotedblright{} (CDIAL 7700), and \emph{abhiṣeka}, \textquotedblleft{}anointing; the inauguration of a king\textquotedblright{} (Monier-Williams).
\end{cousinweb}

\begin{grammarlens}[title={What you've built}]
One more word for telling what happened.
\end{grammarlens}

\subsection*{Your turn}
\begin{itemize}
\raggedright
  \item \emph{Say it:} \emph{paṭṭābhiṣēkaṁ}
  \item \emph{Say it:} \emph{paṭṭābhiṣēkaṁ}, once more
  \item \emph{Say it:} \emph{kotta rāju paṭṭābhiṣēkaṁ ninna jarigindi.} — The new king's coronation took place yesterday.
  \item \emph{Recall:} say the Telugu for a crown, then the Telugu for a throne, then say \emph{paṭṭābhiṣēkaṁ} again
\end{itemize}

\begin{tcolorbox}[breakable,colback=teal!4,colframe=teal!35!black,title={Before you move on}]
What does \te{పట్టాభిషేకం} mean? (\textquotedblleft{}A coronation\textquotedblright{}.) Say it once more.
\end{tcolorbox}
\section[\texorpdfstring{\te{పట్టం} \te{కట్టు} (paṭṭaṁ kaṭṭu)}{పట్టం కట్టు (paṭṭaṁ kaṭṭu)}]{\te{పట్టం} \te{కట్టు} (paṭṭaṁ kaṭṭu) — to enthrone; to crown}

\label{lesson:TE-C335-pattam-kattu}



\begin{quote}
Before the new one: say the Telugu for a throne, then the Telugu for a coronation.
\end{quote}

\subsection*{You'll want to know: \te{పట్టం} \te{కట్టు}}
\textbf{\te{పట్టం} \te{కట్టు}} — \emph{paṭṭaṁ kaṭṭu} — \textquotedblleft{}to enthrone; to crown\textquotedblright{}.

\begin{quote}
\te{ప్రజలు} \te{కొత్త} \te{రాజుకు} \te{పట్టం} \te{కట్టారు}.
\end{quote}

\begin{quote}
\emph{prajalu kotta rājuku paṭṭaṁ kaṭṭāru.}
\end{quote}

\begin{quote}
The people crowned the new king.
\end{quote}

Past (I …): \te{పట్టం} \te{కట్టాను}, \emph{paṭṭaṁ kaṭṭānu}.

Who is crowned takes -\te{కు}: \te{రాజుకు} \te{పట్టం} \te{కట్టు}. It also means to honour a winner.

\begin{cousinweb}
\textbf{\te{పట్టం} \te{కట్టు}} (\emph{paṭṭaṁ kaṭṭu}) is a Telugu phrase: \emph{paṭṭaṁ}, from Sanskrit \emph{paṭṭa}, \textquotedblleft{}a band, a diadem\textquotedblright{} (CDIAL 7700), and native \emph{kaṭṭu}, \textquotedblleft{}to tie\textquotedblright{} (DEDR 1147): to tie the royal band on someone.
\end{cousinweb}

\begin{grammarlens}[title={What you've built}]
One more word for telling what happened.
\end{grammarlens}

\subsection*{Your turn}
\begin{itemize}
\raggedright
  \item \emph{Say it:} \emph{paṭṭaṁ kaṭṭu}
  \item \emph{Say it:} \emph{paṭṭaṁ kaṭṭu}, once more
  \item \emph{Say it:} \emph{prajalu kotta rājuku paṭṭaṁ kaṭṭāru.} — The people crowned the new king.
  \item \emph{Recall:} say the Telugu for a throne, then the Telugu for a coronation, then say \emph{paṭṭaṁ kaṭṭu} again
\end{itemize}

\begin{tcolorbox}[breakable,colback=teal!4,colframe=teal!35!black,title={Before you move on}]
What does \te{పట్టం} \te{కట్టు} mean? (\textquotedblleft{}To enthrone; to crown\textquotedblright{}.) Say it once more.
\end{tcolorbox}
\section[\texorpdfstring{\te{పాలన} (pālana)}{పాలన (pālana)}]{\te{పాలన} (pālana) — a reign; governance}

\label{lesson:TE-C335-palana}



\begin{quote}
Before the new one: say the Telugu for a coronation, then the Telugu for to enthrone.
\end{quote}

\subsection*{You'll want to know: \te{పాలన}}
\textbf{\te{పాలన}} — \emph{pālana} — \textquotedblleft{}a reign; governance\textquotedblright{}.

\begin{quote}
\te{ఆ} \te{రాజు} \te{పాలనలో} \te{ప్రజలు} \te{సుఖంగా} \te{ఉండేవారు}.
\end{quote}

\begin{quote}
\emph{ā rāju pālanalō prajalu sukhaṅgā uṇḍēvāru.}
\end{quote}

\begin{quote}
Under that king's rule the people lived in comfort.
\end{quote}

\te{పాలనలో}, under the rule of: the locative -\te{లో}.

\begin{cousinweb}
\textbf{\te{పాలన}} (\emph{pālana}) is borrowed from Sanskrit \emph{pālana}, \textquotedblleft{}protecting, nourishing\textquotedblright{} (CDIAL 8128): a good ruler protects.
\end{cousinweb}

\begin{grammarlens}[title={What you've built}]
One more word for telling what happened.
\end{grammarlens}

\subsection*{Your turn}
\begin{itemize}
\raggedright
  \item \emph{Say it:} \emph{pālana}
  \item \emph{Say it:} \emph{pālana}, once more
  \item \emph{Say it:} \emph{ā rāju pālanalō prajalu sukhaṅgā uṇḍēvāru.} — Under that king's rule the people lived in comfort.
  \item \emph{Recall:} say the Telugu for a coronation, then the Telugu for to enthrone, then say \emph{pālana} again
\end{itemize}

\begin{tcolorbox}[breakable,colback=teal!4,colframe=teal!35!black,title={Before you move on}]
What does \te{పాలన} mean? (\textquotedblleft{}A reign; governance\textquotedblright{}.) Say it once more.
\end{tcolorbox}
